\section{Auto Research: el comienzo de la investigación automatizada}

\subsection{Idea central: eliminar el cuello de botella humano}

La tesis central de Karpathy: \textbf{para maximizar el aprovechamiento de las herramientas actuales, hay que eliminarse a uno mismo del ciclo}. No se puede quedar ahí esperando para dar el siguiente prompt; hay que dejar que el sistema funcione de manera completamente autónoma. El objetivo es maximizar el rendimiento de tokens y sacar al humano del loop.

\begin{importantbox}{Definición de Auto Research}
Dado un objetivo, una métrica de evaluación y unos límites de operación, se deja que el agent explore de manera autónoma y cíclica posibles mejoras. El humano solo necesita «arrange it once and hit go». La pregunta central es: cómo lograr que más agents se ejecuten durante más tiempo, con menos participación humana.
\end{importantbox}

\subsection{Experimento de DataChat GPT-2}

Karpathy realizó experimentos de auto research con su proyecto DataChat (un playground de entrenamiento de GPT-2). Ya había ajustado el modelo, de manera tradicional (búsqueda manual de hiperparámetros, dos décadas de experiencia), hasta un estado que él mismo consideraba bueno, y luego dejó que el auto researcher se ejecutara durante una noche:

\begin{itemize}
    \item El auto researcher descubrió el weight decay sobre los value embeddings, algo que él había pasado por alto
    \item Los betas de Adam no estaban suficientemente ajustados
    \item Existe una interacción conjunta entre estos hiperparámetros: ajustar uno puede requerir volver a ajustar los demás
\end{itemize}

\begin{warningbox}{El exceso de confianza de los investigadores humanos}
Karpathy señala que incluso un investigador con dos décadas de experiencia en entrenamiento comete errores por descuido. Más importante aún: el humano no debería ser el cuello de botella de la búsqueda de hiperparámetros: «I shouldn't be looking at the results. There's objective criteria.»
\end{warningbox}

\subsection{Program.md: metaoptimización}

La configuración de auto research se describe mediante un archivo \texttt{program.md}, «el ADN de la organización de investigación». Distintos program.md producen distintos avances de investigación.

\begin{importantbox}{Una organización de investigación como código}
«A research organization is a set of markdown files that describe all the roles and how the whole thing connects.» Se puede imaginar que se optimiza el propio program.md: es posible observar al 100\% de dónde vienen las mejoras y luego ajustar el program.md para que se descubran más mejoras de ese tipo. Este es el nivel de la meta-optimization.
\end{importantbox}

\subsection{Límites de aplicación de Auto Research}

\begin{warningbox}{Dos grandes limitaciones de Auto Research}
\begin{enumerate}
    \item \textbf{Se necesita una métrica verificable}: escribir un CUDA kernel (mismo comportamiento, pero más rápido) es el escenario perfecto; pero objetivos blandos como «la comprensión de la intención del usuario» o «el buen gusto del producto» no se prestan al auto research: «if you can't evaluate, you can't auto research it»
    \item \textbf{El modelo todavía es rough around the edges}: si se avanza demasiado lejos y demasiado rápido, el sistema deja de ser útil. El agent desperdicia una gran cantidad de compute en problemas evidentes, lo cual resulta frustrante
\end{enumerate}
\end{warningbox}

\subsection{Resumen de la sección}

Auto research representa un cambio de paradigma: de «el humano investiga manualmente dentro del ciclo» a «el humano configura una sola vez y el agent recorre el ciclo de manera automática». En los dominios con métricas verificables ya ha mostrado resultados sorprendentemente buenos, pero los desafíos de los objetivos blandos y de la jaggedness del modelo todavía deben afrontarse.
